% !TeX root = ../papers/nc-p3.tex
\section{可換な共通二平面での指定階数一行列}
\label{sec:marked-commuting-intersection}
\begin{lemma}\label{lem:three-bracket-rank-one-compression}
$U$ を三次元、$P=u b^t$ を非零階数一行列とする。
既約な跡零三平面 $\mathcal Q\subset\mathfrak{sl}(U)$ の交換子の包が三次元で、
その全元が $W=\ker b^t$ を $L=ku$ に写すなら
\[
                        P_0=P-\operatorname{tr}(P)I/3\ \in\mathcal Q .
\]
\end{lemma}
\begin{proof}
まず $U=W\oplus L$ の場合、全交換子は $2+1$ ブロックの対角二成分が零である。
$X\in\mathcal Q$ を $\left(\begin{smallmatrix}A_X&p_X\\q_X^t&\alpha_X\end{smallmatrix}\right)$ とする。
最後の対角成分より $(p_X,q_X)$ の像は自然な四次元シンプレクティック空間の等方部分なので次元2以下。
非零の $X$ で $p_X=q_X=0$ とできる。全 $Y$ で $[A_X,A_Y]=0$。
$A_X$ が非スカラーならその二次元中心化環は可換であり、全 $A_Y$ が可換となる。
すると左上の式は全 $Y,Z$ に対して $p_Y q_Z^t=p_Z q_Y^t$ である。
$p,q$ の一方が恒等零なら $W$ または $L$ が共通不変部分空間となる。
それ以外では両方非零となる一元を取り、この等式により $(p_Y,q_Y)$ の像が一次元に限られる。
従って二次元以上のブロック対角部分があり、そこで全元が可換なので交換子の次元3に反する。
よって $A_X$ はスカラーであり、跡零の非零 $X$ が $P_0$ の倍である。

次に $L\subset W$ とし、$P=E_{13}$ に取る。$\mathcal Q$ のパラメータ空間上の一次式を用いて行列を
$\left(\begin{smallmatrix}a&b&c\\d&e&f\\g&h&i\end{smallmatrix}\right)$ と書く。
交換子の左下二行二列（第2,3行、第1,2列）の零は
\[
 d\wedge(a-e)+f\wedge g=0,\quad
 g\wedge(a-i)+h\wedge d=0,\quad
 g\wedge b+h\wedge(e-i)=0,\quad d\wedge b+f\wedge h=0 .
\]
$g\ne0$ なら二番目の式から $d,g,h$ の包の次元は高々2。
次元1なら $d,h$ は $g$ の倍であり、旗を保つ基底変更で $d=h=0$ とする。
実際 $I+tE_{23}$ による共役は $d$ に $tg$ を加え、$I+t'E_{12}$ による共役は
$h$ から $t'g$ を引くのでこの変更ができる。その基底では四式より $f,b,a-i$ も $g$ の倍なので $\ker g$ の二平面は
$\operatorname{diag}(a,e,a)+cE_{13}$ の形で可換となり矛盾する。
従って次元2である。$d,g,h$ を全て消す非零のパラメータ $s$ で四式を縮約すると
$a_s=i_s$, $f_sg=(a_s-e_s)d$, $b_sg=(a_s-e_s)h$ となる。
二次元性から $a_s=e_s$、さらに $f_s=b_s=0$。従って $Q_s$ は $E_{13}$ の非零倍である。
$g=0$ なら $h\wedge d=0$ で、既約性より両方非零だから $h=kd$, $k\ne0$。
残りの式より $a-e,e-i,b-kf$ は $d$ の倍である。$\ker d$ の二平面は
$kE_{12}+E_{23}$ と $E_{13}$ の包となって可換であり、再び交換子の次元に反する。
\end{proof}

\begin{theorem}\label{thm:marked-commuting-intersection}
命題\ref{prop:marked-mixed-discriminant}の三次式が恒等零のとき、
$H=G\cap F$ に対して $W=\{c:F_c\in H\}$、$L=\{a:G_a\in H\}^\perp$ とする。
$L$ を像、$W$ を核とする階数一行列 $P$ は定数倍を除いて指定される。
$\langle Q\rangle$ は $P_0$ を含み、その方向 $[s]$ は次の\emph{既知の射影直線}上にある。
指定交換子写像を $B:\bigwedge^2V\longrightarrow\mathcal C$ とする。
$\operatorname{im}[P,-]$ 上の交代形式を
\[
      \Omega([P,Y],[P,Z])=\operatorname{tr}(P[Y,Z])
\]
とし、$\mathcal C$ への制限の根基直線を $D$ とする。
求める直線は $B^{-1}(D)$ の非零二ベクトルの張る二平面の射影化である。
さらに $s$ は $H$ の共通固有方向を取り得ない。
\end{theorem}
\begin{proof}
$g=4$ なので既出の二交換子の次元判定から $B$ は三次元間の同型である。
$a\in L^\perp,c\in W$ に対し $[G_a,F_c]=0$ なので、全ての指定交換子は $W$ を $L$ に写す。
跡零でこの条件を持つ空間は四次元であり、$[P,-]$ の像と等しい。
実際この像はその中にあり、$P$ の中心化環は $P$ が半単純でも平方零でも五次元である
（それぞれ上の二つの基底で読める）。補題より $Q_s$ は $P_0$ の非零倍に選べる。
$\Omega$ は跡の巡回性により代表によらず、$\operatorname{tr}([P,Y] Z)$ の完全対合から
非退化である。$B(s\wedge V)$ 上では $\operatorname{tr}(P[Q_t,Q_u])=0$ のため
$\Omega=0$（全指定値は $P$ と跡直交する）。四次元シンプレクティック空間内の超平面
$\mathcal C$ の根基は直線であり、その中の等方二平面は必ずこの根基を含む。
$B$ で戻せば $s$ の位置の主張を得る。

最後に $Q_s$ が $W$ 上に作用するスカラーを $\lambda$ とする。
もし $s$ が $H$ の共通固有方向なら、$c\in W$ ごとに $F_c s=\mu(c)s$ で
$\mu$ は一次式である。辺 $12$ の二つの点を $(-\lambda,s)$ と $(-\mu(c),c)$ に取ると
六つの独立な式は $(Q_s-\lambda I)c=0$, $(F_c-\mu(c)I)s=0$ として全て成立する。
これは定理\ref{thm:marked-point-commutators}の一方に星印のある六行の式であり、
$F_c$ が対称なので左右は同じである。$[c]\in\mathbf P(W)$ が動くため辺の有限性に反する。
\end{proof}

\begin{theorem}\label{thm:no-three-bracket-compression}
既約な跡零三平面の三つの交換子が独立なら、その全てが固定二平面を固定直線に送ることはない。
従って前定理の可換な共通二平面の場合は有限窓に存在しない。
さらに $q=3$, $g>3$ で交換子が三独立なら、$g=5$ でも判別三次式は恒等零にならない。
\end{theorem}
\begin{proof}[階数一行列を含んだ後の排除]
補題により $P_0\in\mathcal Q$。まず $W\oplus L$ の場合に、
その補基底 $Y,Z$ の左上二行二列を $A,B$、右上列を $p,r$、左下行を $q^t,t^t$ とする。
$[Y,Z]$ の対角の消失より
\[
               q^t r=t^t p,\qquad [A,B]=r q^t-p t^t .
\]
$[P,Y],[P,Z]$ は独立であり、$(p,q),(r,t)$ は二平面を張る。
$p,r$ が独立なら $q=Sp,t=Sr$ と書け、第一式で $S$ は対称である。
$S=0$ なら $W$ が不変になる。$J=r p^t-p r^t$ は可逆な交代行列で $[A,B]=JS$。
$A,B$ はそれぞれその交換子と跡直交するので、$SA,SB$ は対称である。
$S$ が特異なら $\ker S\subset W$ は共通不変直線となる。
可逆なら $P_0,Y,Z$ は計量 $\operatorname{diag}(S,1)$ に関して自己随伴である。
しかしこの計量に関する交代行列で両対角ブロックが零のものは二次元しかなく、
三交換子の独立性に反する。$q,t$ が独立なら転置して同じである。
どちらの射影も階数1なら、二基底を取り直して
$(p,q)=(p,0)$, $(r,t)=(0,t)$, $p,t\ne0$ とできる。$t^tp=0$ で
$[A,B]=-pt^t$。跡直交により $A,B$ は $\ker t^t=kp$ を保ち、
三行列全てがこの直線を保つ。射影が零の場合も一方のブロックが不変である。

$P=E_{13}$ の場合は補題の四式の記号を使う。既に $g\ne0$ かつ
$d,g,h$ の包は二次元である。$g_Y=1$, $g_Z=0$ として $P_0,Y,Z$ を基底に取る。
旗を保つ先の二つの共役で $d_Y=h_Y=0$ とし、$I+vE_{13}$ による共役で
$a_Y=i_Y$ にする。この最後の変更は $d_Y,h_Y$ を保ち、$a_Y-i_Y$ に $2v$ を加える。
$D=d_Z$, $H=h_Z$ は同時には零でない。四式を $Y,Z$ で評価すると
\[
 a_Z=i_Z,\qquad f_Z=D(e_Y-a_Y),\qquad b_Z=H(e_Y-a_Y),
 \qquad D b_Y=H f_Y .
\]
$D=0$ なら $d=f=0$ で $\langle e_1,e_3\rangle$ が不変、$H=0$ なら $h=b=0$ で
$ke_2$ が不変である。両方非零なら $j=H/D$ として三行列とも
\[
                  M=\begin{pmatrix}0&0&1\\0&j&0\\1&0&0\end{pmatrix}
\]
に関して自己随伴である。実際 $h=jd$, $b=jf$, $i=a$ が必要十分で、$c$ は自由である。
一方、左下二行二列が零の跡零行列で $M$ に関して交代となるものは二次元である。
これで両方の配置で三つの指定値は実現できない。
\end{proof}

\begin{proof}[五次元和を含む判別式]
三独立な指定値の空間 $\mathcal C$ に共通右核 $c$ があれば $F_c$ が全 $G_a$ と交換し、
既約性と跡零性で零になり $F$ の独立性に反する。共通左核もない。
三次式が恒等零なら、このような三次元特異行列空間は固定二平面を固定直線に送るものか、
可逆な左右変換で交代三行列の空間になるものである。ここではこの二択を直接確かめる。

階数2の元を左右変換で $A=\operatorname{diag}(1,1,0)$ とし、補二変数を $x$ とする。
全元が階数1以下の場合は共通核または余核があるため除いた。行列式の $A$ での一次式により
補元は $\left(\begin{smallmatrix}H(x)&p(x)\\q(x)^t&0\end{smallmatrix}\right)$ である。
$\det(tA+M(x))=0$ から $q^tp=q^t\operatorname{adj}(H)p=0$。
$p$ が階数1なら $p=w\alpha$、$q$ は $w$ を消す固定行の倍となる。両者とも恒等零にはなれない。
$w=e_1$ とすれば第二式から $H_{21}=0$ で、全元が $\langle e_1,e_3\rangle$ を $ke_1$ に写す。
$p$ が階数2ならそれを $x$ 自身に取り、$q$ は固定可逆交代行列 $J$ で $Jp$ に書ける。
第二式から $Hp=\lambda(x)p$ であり、二成分が互いに素なので $\lambda$ は一次式である。
$H-\lambda I$ の各行は $p$ を消す一次式の行だから $H=\lambda I+u q^t$、$u$ は定数列である。
補元から $\lambda A$ を引き、最終行の倍を上二行から引けば $H=0$。
上二行に $J$、最終行に $-1$ を掛けると全行列が交代となる。

第一の型は定理前半で生じない。第二の型なら元の空間は $\{KB:K^t=-K\}$、
$B$ 可逆と書ける。全行列の跡が零のため $B$ は対称で、$U$ の基底変更により
$\mathcal C$ を交代空間にできる。この基底変更は $G,F$ の対称行列をそれぞれ線型結合で取り直す。
指定値が $a,c$ で交代であることから $[G_x,F_x]=0$ が恒等的に成立する。
補題\ref{lem:commuting-linear-symmetric-nets}により $F_x$ は $G_x$ の一定倍で $g=3$ に反する。
以上で判別式が恒等零となる場合を全て除いた。
\end{proof}
